% Shared exercise chunk: l2-minimax-regret-q
\begin{exercisechunk}
\item \textbf{Expressing simultaneous optimality by minimax regret.}
\label[exercise]{ex:minimax-regret-q}
Consider the excess risk relative to the $Q$-specific minimax value:
\begin{equation}
\min_{\cA}\max_{Q\in\cQ}
\left\{
\max_f R_Q(\cA,f)-V(Q)
\right\}.
\label{eq:minimax-regret-q}
\end{equation}
\begin{enumerate}[(a)]
\item Show that the value of \cref{eq:minimax-regret-q} is zero.
\item Prove that a learner
attains this value exactly when it is minimax optimal simultaneously for every
$Q$.
\item Using \cref{eq:q-specific-value}, characterize its prediction at every
seen and unseen test index for every consistent labeled sample.

\item Compute the worst-case regret of the learner that always predicts
with a fair coin and determine whether it minimizes \cref{eq:minimax-regret-q}.
Compare its optimality here with its optimality for
\cref{eq:absolute-distribution-free}.
\end{enumerate}
\end{exercisechunk}
